\section{GB9: The Lean-Checked Core and the Scope-Boundary Theorem}
\label{sec:lean}

\subsection{What is machine-checked}

Lean \cite{de-moura-ullrich-2021} is a proof assistant: a theorem stated in it is accepted only
when a machine has checked every step, and a development ``compiles sorry-free'' when no step has
been left as an unproved placeholder. We use it for two things. The first is the abstract theorems
the certificates rely on, for example that a witness pair exists if and only if the predictive
description is strictly finer than the present one. The second is the logical consequences of the
exact witness tables that the computation exports. Neither then depends on a reader trusting the
prose. Precisely:

\claimref{GB9}{712fa267299cef1fce505d5b88bb7689edac6f3cadbcdc932e371592a24ce78a}
Three Lean developments compile sorry-free on the pinned toolchain,
\texttt{leanprover/\allowbreak{}lean4:4.29.0}: the vendored legacy \texttt{HolonomyMemory} core (paper [A]'s
development \sbtcite{sbt-paper-a}); the \texttt{DeclaredMemory} core, adapted from the upstream
route-transport development by renaming modules and namespace only, with its mathematical
declarations unchanged; and this project's \texttt{GlassWitness} extension. All $96$ theorems carry
per-theorem axiom reports, regenerated from the live build, that match the stored audit. No
theorem depends on anything beyond the standard axioms \texttt{propext},
\texttt{Classical.choice} and \texttt{Quot.sound}, and no \texttt{native\_decide}, \texttt{sorry}
or custom axiom appears. Per the framework's fidelity legend, each row of the $97$-row fidelity
table (Appendix~\ref{app:lean}; $96$ theorems plus one row for evidence that lives outside Lean)
carries its true label: $70$ \emph{verified}, $25$ \emph{narrowed/parametric} (the finite rational
lookup witnesses and their generic constructors), $1$ \emph{schema projection}, and $1$
\emph{computational evidence only}. No schema or trivial declaration is described as a verified
proof, and closure-deficit computations are not Lean theorems.

Among the verified rows are the witness/strict-refinement equivalence
(\texttt{strictRefinement\_\allowbreak{}iff\_\allowbreak{}nonempty\_\allowbreak{}predictiveWitness}) and the factorization theorems for
future-sufficient states. Their direction matters: the minimal predictive quotient $M$ factors,
uniquely on the reachable image, through \emph{any} future-sufficient state $S$. The converse need
not hold, since $S$ may carry redundant information.

\subsection{The witness tables}

The numeric witnesses are checked in the \texttt{DeclaredMemory} core as finite tables of exact
rationals, exported by a generator from the Python computation:

\begin{itemize}
  \item \texttt{NumericN10}: the GT2 Kovacs pair, current mean energy $20/7$ and the two exact
    five-step future values.
  \item \texttt{NumericN8}: a mean-energy pair of the same kind on the $N=8$ GT6 shell
    configuration (current mean energy $16/7$).
  \item \texttt{EnergyN8}: the GT6 pair, with nine current events (the complete energy
    distribution, shared by both histories) and $21$ future horizons.
\end{itemize}

For each, Lean checks the rational inequalities by kernel reduction and proves strict refinement,
nonfactorization of the future through the current readout, and failure of future sufficiency
for the current readout. For \texttt{EnergyN8} this covers every horizon and every descriptor
computed from the current energy law. Generic constructors prove the same consequences for any
table with an actual rational split, and a constant-future negative control shows that a table
without a split yields no strictness.

\claimref{NEG.lean\_instance\_structural}{712fa267299cef1fce505d5b88bb7689edac6f3cadbcdc932e371592a24ce78a}
The boundary of this mechanization is precise. Lean checks the exported tables and what follows
from them. It does not compute the East kernel, the preparations, the randomized-stop tuning, the
null-vector selector, or the equality of the two current energy laws in GT6; those are verified in
exact Python before export, and the fidelity table carries them as its one
computational-evidence row. The legacy \texttt{glassInstance} remains a typed, index-based
structural mirror.

\subsection{The scope-boundary theorem}

Why a second core? Because in the legacy core no witness can exist at all, and Lean proves it.

\claimref{FINDING.lean\_collapse}{b29d612b26eddaa25183177e66a480e4b317ebd3e7794e16f1776a1e8ac08e11}
For \emph{every} instance of the legacy \texttt{RouteTransportCore}, the \texttt{observe\_push}
coherence field, combined with \texttt{CurrentEventEquiv}'s universal quantification over events,
forces current equivalence to imply future-predictive equivalence
(\texttt{current\_\allowbreak{}implies\_\allowbreak{}future\_\allowbreak{}GENERIC}). Consequently the map from predictive to current
classes is injective, and \texttt{PredictiveWitness}, \texttt{StrictRefinement} and
\texttt{LoopAsymmetry} instances cannot be constructed in that core
(\texttt{predictiveWitness\_\allowbreak{}impossible\_\allowbreak{}GENERIC}, \texttt{strictRefinement\_\allowbreak{}impossible\_\allowbreak{}GENERIC},
\texttt{predictiveToCurrent\_\allowbreak{}injective\_\allowbreak{}GENERIC}, \texttt{loopAsymmetry\_\allowbreak{}impossible\_\allowbreak{}GENERIC};
machine-checked, sorry-free).

The mechanism is simple. The legacy core's naturality axiom lets every future event be pulled back
to a present event, so ``agreement on all present events'' silently includes agreement on all
pulled-back futures, and nothing is left for a witness to separate. The \texttt{DeclaredMemory}
core drops the pullback and its coherence field: future observations are declared events of the
pushed history rather than present events in disguise. That single change is what admits the
tables above. The Python certificates never had the problem, because they quantify over a finite
\emph{declared} catalog of futures rather than over all pullback-closed events.

The collapse has a familiar shape. It is the proof pattern of strong lumpability for Markov chains
\cite{kemeny-snell-1960,buchholz-1994}, where a local one-step coherence condition propagates an
equivalence to all future times. It also has the form of biextensional collapse in the Chu-space
sense \cite{barr-1979,pratt-1999,parente-2024}, with \texttt{pullEvent} playing a role analogous to
the Chu transform's naturality condition; we offer this reading as an analogy, and no formal
correspondence is claimed or proved. The instructive contrast is the causal-state construction of
computational mechanics \cite{shalizi-crutchfield-2001}, where agreement on the present generally
does \emph{not} imply agreement on the future: the collapse is forced by the legacy core's axiom,
not by prediction itself.
